\chapter{Age-Related Macular Degeneration}
\index{Age-Related Macular Degeneration}
\label{sec:Age-Related Macular Degeneration}

\section{Overview}
This degenerative condition involves progressive deterioration of affected tissues, leading to gradual loss of function. It is more common with advancing age.
\section{Causes}
This condition happens because of age-related degenerative changes in the retinal pigment epithelium and photoreceptors. Degenerative processes involve progressive cellular dysfunction and tissue breakdown over time. Contributing factors include oxidative stress, inflammation, accumulated cellular damage, and reduced regenerative capacity.    
\section{Risk Factors}

\begin{itemize}[noitemsep]
  \item Advancing age
  \item Genetic predisposition
  \item Repetitive use or overuse (joints)
  \item Previous injuries
  \item Obesity
  \item Smoking and alcohol use
  \item Genetic predisposition and family history
  \item Lifestyle factors (diet, exercise, smoking, alcohol)
  \item Environmental exposures
  \item Underlying medical conditions
  \item Medication use
\end{itemize}
\section{Signs and Symptoms}

\subsection{Common symptoms}
\begin{itemize}[noitemsep]
  \item You may experience progressive central vision loss and distortion. 
  \end{itemize}

\subsection{Emergency warning signs}
\begin{itemize}[noitemsep]
  \item Severe or worsening symptoms of age-related macular degeneration require immediate medical evaluation
\end{itemize}
\section{Progression and Stages}
The condition may progress through different stages of severity over time. Early stages often involve mild or subtle symptoms that may be overlooked. Moderate stages involve more noticeable symptoms affecting daily function. Advanced stages may involve significant impairment and complications.

\begin{itemize}[noitemsep]
  \item Age-related macular deterioration is the leading cause of irreversible central vision loss in individuals over 50 in developed countries, affecting the macula. Early AMD is marked by drusen and pigmentary changes
  \item Advanced AMD manifests as either dry (atrophic) AMD with geographic shrinkage or wasting or wet (neovascular) AMD with choroidal neovascularization.
\end{itemize}
\section{Complications}
\begin{itemize}[noitemsep]
  \item Disease progression leading to worsening symptoms
  \item Secondary infections or injuries
  \item Side effects of treatment
  \item Reduced quality of life
  \item Emotional and psychological impact
\end{itemize}
\section{Diagnosis}
\begin{itemize}[noitemsep]
  \item Doctors diagnose this based on fundus examination and optical coherence tomography. Treatment for wet AMD involves anti-VEGF intravitreal injections (ranibizumab, aflibercept, bevacizumab)
  \item Complement inhibitors are under investigation for dry AMD.
  \item Comprehensive medical history and physical examination
  \item Laboratory tests to assess organ function
  \item Imaging studies to visualize affected structures
  \item Specialized testing depending on the affected system
  \item Consultation with relevant specialists
\end{itemize}
\section{Prevention}
\begin{itemize}[noitemsep]
  \item Maintain a healthy lifestyle with balanced diet and exercise
  \item Avoid known risk factors
  \item Attend regular medical check-ups for early detection
  \item Follow recommended screening guidelines
\end{itemize}
\section{Treatment Options}
Medications to manage symptoms and address underlying mechanisms Lifestyle modifications and self-care measures Physical therapy and rehabilitation Surgical intervention when indicated Regular monitoring and follow-up care Patient education and support services

\begin{itemize}[noitemsep]
  \item Medications to manage symptoms and address underlying mechanisms
  \item Lifestyle modifications and self-care measures
  \item Physical therapy and rehabilitation
  \item Surgical intervention when indicated
  \item Regular monitoring and follow-up care
\end{itemize}
\section{Living with It}
\begin{itemize}[noitemsep]
  \item Follow your treatment plan as prescribed
  \item Maintain regular follow-up with your healthcare provider
  \item Make necessary lifestyle adjustments
  \item Seek support from family, friends, or support groups
  \item Stay informed about your condition
\end{itemize}
\section{Outlook}
The outlook has improved significantly with anti-VEGF therapy for wet AMD. The outlook varies depending on the specific condition, severity at diagnosis, and response to treatment. Early intervention generally leads to better outcomes. Many conditions can be effectively managed with appropriate treatment. Regular follow-up care is important for monitoring and treatment adjustment.
